% 原书分章 PDF 第 21—24 页：第 25.6 节，完整密度估计例子与定理 25.5。
\section{例：一元密度的单点估计}

令 $O\sim P_0$，$\mathcal M$ 为非参数模型，
对于给定点 $x_0$，$\Psi(P_0)=p_0(x_0)$。
定义 $\Psi_h(P)=\int_tK_{h,x_0}(t)p(t)\,dt$，
其中 $K$ 为 $J-1$ 阶正交核，
$K_{h,x_0}(t)=h^{-1}K((t-x_0)/h)$。
$\Psi_h$ 在 $P$ 处的有效影响曲线为
$D_h^*(P)=K_{h,x_0}(O)-\Psi_h(P)$。
它满足恒等式
$\Psi_h(P)-\Psi_h(P_0)=-P_0D_h^*(P)$，
因此 $R_{20,h}(P,P_0)=0$。
令 $\mathbf p_n$ 为 $p_0$ 的初始估计量。
令 $B_n\in\{0,1\}^n$ 为交叉验证方案，
$\mathbf p_{n,B_n}$ 为将初始密度估计量应用于训练样本
$\{O_i:B_n(i)=0\}$ 所得的估计。
考虑在 $\epsilon=0$ 时经过 $p$ 的局部最不利子模型：
\[
 \mathbf p_{\epsilon,h}(o)
 =(1+\epsilon D_h^*(P)(o))p(o).
\]
可以使用这个局部最不利子模型，也可以使用通用最不利子模型
\[
 \mathbf p_{\epsilon,h}(o)=p(o)
 \exp\left(\int_0^\epsilon D_h^*(p_e)(o)\,de\right).
\]
令 $\mathbf p_{n,B_n,\epsilon,h}$ 为经过
$\mathbf p_{n,B_n}$ 的这个最不利子模型。令
\[
 \epsilon_n=\argmin_\epsilon
 E_{B_n}P_{n,B_n}L(\mathbf p_{n,B_n,\epsilon,h}),
\]
其中 $L(p)(o)=-\log p(o)$ 为对数似然损失。
这定义了 $\mathbf p^*_{n,B_n,h}=\mathbf p_{n,B_n,\epsilon_n,h}$，
从而得到 $\Psi_h(P_0)$ 的 CV-TMLE
$\psi^*_{h,n}=E_{B_n}\Psi_h(\mathbf P^*_{n,B_n})$。
将 $h_n$ 定义为下式的一个解：
\begin{equation}
 \frac{d}{dh_n}\psi_{h_n,n}
 +C_n\sigma_n\frac{d}{dh_n}(nh_n)^{-0.5}=0,
\tag{25.11}
\end{equation}
其中将 $\psi_{h,n}$ 定义为 CV 一步估计量：
\[
 \psi_{h,n}=E_{B_n}P^1_{n,B_n}K_{h,x_0}
 =P_nK_{h,x_0}
 =\frac{1}{nh}\sum_{i=1}^{n}K((x_0-O_i)/h).
\]
按照我们对该方法的描述，对于某个序列 $\alpha_n$，可设
$C_n\in\{-Z_{1-\alpha_n},Z_{1-\alpha_n}\}$。
此外，$\sigma_n^2=\tilde p_n(x_0)\int K(y)^2\,dy$
是 $\sigma_0^2=p_0(x_0)\int K(y)^2\,dy$ 的相合估计量。
这里，$\tilde p_n(x_0)$ 可以替换为任意相合估计量，
例如采用某个满足 $nh_n^0\to\infty$ 的带宽选择器 $h_n^0$
得到的 CV-TMLE $\psi^*_{h_n^0,n}$。
注意，由于 $K_{h,x_0}$ 不依赖在训练样本上估计的干扰参数，
带宽选择器 $h_n$ 中使用的 CV 一步估计量，
简化为 $x_0$ 处的标准核密度估计量（kernel density estimator）。
我们所提议的 CV-TMLE 为 $\psi_n^*=\psi^*_{h_n,n}$。

现在应用定理 25.1。
假定 $p_0(x_0)>0$，从而在 $t_0$ 的一个小邻域内也远离零。
先考虑条件 (25.3)。有
\[
\begin{aligned}
 \|D_h^*(\mathbf P_n)-D_h^*(P_0)\|_{P_0}
 &=|\Psi_h(\mathbf P_n)-\Psi_h(P_0)|\\
 &=\left|\int_xh^{-1}K((x-x_0)/h)
 (\mathbf p_n(x)-p_0(x))\,dx\right|\\
 &=\left|\int K(y)
 (\mathbf p_n(x_0+hy)-p_0(x_0+hy))\,dy\right|\\
 &=O\left(h^{-1}\int_{x_0-h}^{x_0+h}
 |(\mathbf p_n-\mathbf p_0)(x)|\,dx\right).
\end{aligned}
\]
例如，若 $\mathbf p_n-\mathbf p_0$ 在 0 周围的一个小区间上，
以速率 $r_1(n)$ 一致趋于零，则右端为 $O_P(r_1(n))$。
也可以利用 Cauchy--Schwarz 不等式，作更保守的控制：
\[
 h^{-1}\int I_{(x_0-h,x_0+h)}|\mathbf p_n-p_0|\,dx
 \leq O\left(h^{-0.5}
 \left(\int(\mathbf p_n-p_0)^2\,dx\right)^{0.5}\right),
\]
其中积分取遍一个包含 $x_0$ 的固定区间。
此时，只需要对于某个多项式速率，
$\|\mathbf p_n-p_0\|_{t_0,P_0}=O_P(r(n))$，
就能得到界 $O_P(h^{-0.5}r(n))$。
这里，对固定的 $a>0$，定义
$\|p-p_0\|_{x_0,P_0}^2
=\int_{x_0-a}^{x_0+a}(p-p_0)^2\,dP_0$。
这验证了第一个条件 (25.3)。注意，
\[
\begin{aligned}
 hP_0D_h^*(p_0)^2
 &=hP_0\{K_{h,x_0}-\Psi_h(P_0)\}^2\\
 &=hP_0K_{h,x_0}^2+O(h)\\
 &=\int K(y)^2p_0(x_0+hy)\,dy+O(h),
\end{aligned}
\]
它收敛到 $\sigma_0^2=p_0(x_0)\int K(y)^2\,dy$。
由于 $R_{20,h}=0$，条件 (25.5) 自动成立。
因此，现在可以应用定理 25.1，证明
\[
 \sqrt{nh}(\psi^*_{h,n}-\psi_{h,0})
 =\sqrt{nh}(P_n-P_0)D_h^*(P_0)+o_P(1)
 \Rightarrow_dN(0,\sigma_0^2).
\]
假定 $p_0$ 在 $x_0$ 处 $J_0$ 次连续可微，并令 $J_0^*=\min(J_0,J)$。
则有 $h^{-J_0^*}(\psi_{h0}-\psi_0)\to B_0(J_0^*)$，其中
\begin{equation}
 B_0(J_0^*)\equiv
 \psi_0^{J_0^*}(x_0)\int_yK(y)\frac{y^{J_0^*}}{J_0^*!}\,dy.
\tag{25.12}
\end{equation}
应用定理 25.2 可证明，若选择最优速率
$h_n^*=n^{-1/(2J_0^*+1)}$，则
\[
 n^{-J_0^*/(2J_0^*+1)}
 (\psi^*_{h_n,n}-\psi_0)\Rightarrow_dN(B_0(J_0^*),\sigma_0^2).
\]
由于定理 25.2 的条件与定理 25.1 相同，
只是把 $K_{h,x_0}$ 换成
$K_{1,h,x_0}=h\frac{d}{dh}K_{h,x_0}$，
而我们已证明后者是行为类似的核，因此条件的验证与上面相同。
应用定理 25.3 可证明
$h_n\sim(C_n^2/n)^{-1/(2J_0^*+1)}$；
若 $C_n=O_P(1)$，它就达到最优速率。这证明了以下定理。

\begin{bookresult}{定理}{25.5}
考虑上面定义的 CV-TMLE $\psi^*_{h,n}$，
以及式 (25.11) 的带宽选择器 $h_n$。
并考虑 $B_0(J)$ 的定义 (25.12) 与 $\sigma_0^2$。作如下假设。
\begin{itemize}
\item $p_0(x_0)>0$。
\item 对某个未知的 $J_0$，$p_0$ 在 $x_0$ 处 $J_0$ 次连续可微。
\item 对某个多项式速率，即某个 $p>0$ 对应的 $n^{-p}$，有
$\|\mathbf p_n-p_0\|_{x_0,P_0}=O_P(r(n))$；
其中对固定的 $a>0$，
$\|p-p_0\|_{x_0,P_0}^2
=\int_{x_0-a}^{x_0+a}(p-p_0)^2\,dP_0$。
\end{itemize}
令 $J_0^*=\min(J_0,J)$。有
$h_n\sim(C_n^2/n)^{-1/(2J_0^*+1)}$；
并且若 $C_n=O_P(1)$，则
\[
 (nh_n^d)^{0.5}(\psi^*_{h_n,n}-\psi_0)
 \Rightarrow_dN(B_0(J_0^*),\sigma_0^2).
\]
假定 $C_n=O_P(1)$，并令 $h_n^*=h_n/\log n$
为 $h_n$ 的欠平滑版本。
令 $\psi^*_{h_n^*,n}$ 为相应的 CV-TMLE。则
\[
 (nh_n^*)^{-0.5}(\psi^*_{h_n^*,n}-\psi_0)
 \Rightarrow N(0,\sigma_0^2).
\]
特别地，
\[
 \psi^*_{h_n^*,n}\pm1.96\sigma_n/(nh_n^*)^{0.5}
\]
是 $\psi_0$ 的渐近 $95\%$ 置信区间。
\end{bookresult}
